\section{Conclusions}
\label{sec:conclusions}
In our project, we attempted to find a difficulty metric for Wikispeedia games. We found reasonable success, as long as the metric takes into account a combination of graph structure and semantic distances as given by our LLM embeddings. Metrics that look at these factors in isolation performed worse as a predictor of difficulty.


% % ### Intuititions for future directions I got from ChatGP
% % I think we did not select a good noise, the ones outlined below are some alternatives we can briefly suggest in future work. Especially (a) and (d) make sense to me.
%
% \subsection*{(a) Semantic Reach Centrality (Global Semantic Closeness)}
%
% We define the effective width of a target node $t$ as
% \begin{equation}
% W_{\text{target}} = \frac{1}{N} \sum_{i \in V} \exp\left(-\lambda \, d_{\text{sem}}(i, t)\right),
% \end{equation}
% where $d_{\text{sem}}(i, t)$ is the shortest-path semantic distance between node $i$ and the target $t$,
% computed using link weights based on semantic similarity.
%
% \noindent \textbf{Intuition:}
% A target with many semantically close nodes has a large width, while distant or obscure topics
% have small width. The exponential term ensures diminishing influence with increasing distance.
% This measure is essentially a form of \textit{semantic closeness centrality}, analogous to Katz centrality.
%
% \subsection*{(b) Semantic Katz or Personalized PageRank}
%
% Instead of a uniform random-walk PageRank, we can define:
% \begin{equation}
% W_{\text{target}} = \mathrm{KatzCentrality}_{\text{sem}}(t)
% = \sum_{k=1}^{\infty} \beta^k
%     \sum_{\substack{\text{paths of length } k \\ \text{ending at } t}}
%     \prod_{e \in \text{path}} s(e),
% \end{equation}
% where $s(e)$ denotes the semantic similarity on edge $e$, and $\beta \in (0,1)$
% is an attenuation factor. This expression counts semantically weighted paths
% leading to the target --- a natural analog of the ``width of the attractor basin.''
%
% \subsection*{(c) Semantic Basin Size (Thresholded Reachability)}
%
% For a semantic radius $\delta$, define:
% \begin{equation}
% W_{\text{target}} =
% \left|
% \left\{ i \in V : d_{\text{sem}}(i, t) < \delta \right\}
% \right|.
% \end{equation}
% That is, $W_{\text{target}}$ represents how many nodes are semantically within a
% radius $\delta$ of the target. This measure corresponds to the literal size of the target's
% \textit{semantic capture region}.
%
% \subsection*{(d) Empirical Reachability}
%
% If Wikipedia game traversal data are available, $W_{\text{target}}$ can be estimated empirically as
% \begin{equation}
% W_{\text{target}} =
% \frac{1}{N_{\text{start}}}
% \sum_{i}
% \frac{1}{
%     \mathbb{E}\!\left[ \text{path length from } i \text{ to target} \right]
% }.
% \end{equation}
% Targets that are reached quickly from many starting points have larger effective widths.
% This is a behavioral rather than structural measure, but it directly captures
% the observed ease of reaching the target.
